% ------------------------------------------------------------
% §3  The construction and its finite geometry
% ------------------------------------------------------------
\section{Anatomy: a self-map of a finite space}\label{sec:geometry}

We now open the organism. The dissection has two stages, and only the
second is specific to \SHA{}: first a completely general---and genuinely
theoretical---reduction of ``hash anything'' to ``iterate one map on a
finite space''; then the finite space itself and the map that \SHA{}
iterates on it.

\subsection{From arbitrary bit strings to one finite map}
\label{sec:merkledamgard}

A function on $\bigcup_{\ell < 2^{64}} \bitstrings^\ell$ looks
unwieldy---its domain is (for practical purposes) infinite and
ragged. The first structural fact about \SHA{} is that all of this
unwieldiness is handled by a single, provably safe reduction, the
\emph{Merkle--Damg{\aa}rd construction}, and that everything interesting
then happens on a finite set.

\paragraph{Padding.} A message $m$ of bit length $\ell$ (any $\ell$,
including $0$; bits, not bytes) is padded by appending a single
$\mathsf{1}$ bit, then the minimum number of $\mathsf{0}$ bits, then the
number $\ell$ written as a $64$-bit integer, so that the total length is a
multiple of $512$. The padded message splits into $512$-bit blocks
$m_1, \dots, m_k$ (\cref{fig:padding}). Padding is where ``arbitrary bit
length'' lives: the map
\[
  \mathrm{pad} : \bigcup_{\ell < 2^{64}} \bitstrings^{\ell}
  \;\longrightarrow\; \bigl(\bitstrings^{512}\bigr)^{+}
\]
is injective (the appended $\mathsf{1}$ and the length field between them
guarantee it), and encoding the \emph{length} in the final block---so-called
Merkle--Damg{\aa}rd strengthening---is not bookkeeping but the hypothesis
that makes \cref{thm:md} below true.

\begin{figure}[ht]
  \centering
  \begin{tikzpicture}[font=\small, >={Stealth[length=1.8mm]}]
    % ---- the linear message ----
    \draw[fill=msgfill] (0,2.2) rectangle (9.03,2.8);
    \draw (3.7,2.2) -- (3.7,2.8);
    \draw (7.4,2.2) -- (7.4,2.8);
    \node at (1.85,2.5)  {\textsf{A}};
    \node at (5.55,2.5)  {\textsf{B}};
    \node at (8.21,2.5)  {\textsf{C}};
    \node[anchor=west] at (9.25,2.5) {message $m$, $\ell$ bits};
    % ---- guide lines: where each portion lands ----
    \draw[dashed, gray] (0,2.2)    -- (0,0.62);
    \draw[dashed, gray] (3.7,2.2)  -- (3.95,0.62);
    \draw[dashed, gray] (7.4,2.2)  -- (7.9,0.62);
    \draw[dashed, gray] (9.03,2.2) -- (9.53,0.62);
    % ---- the 512-bit blocks ----
    \draw[fill=msgfill] (0,0)    rectangle (3.7,0.6);
    \draw[fill=msgfill] (3.95,0) rectangle (7.65,0.6);
    \node at (1.85,0.3) {\textsf{A}};
    \node at (5.8,0.3)  {\textsf{B}};
    % final block: tail of m | pad bit 1 | 0...0 | length
    \draw[fill=msgfill]  (7.9,0)   rectangle (9.25,0.6);
    \draw[fill=padfill]  (9.25,0)  rectangle (9.75,0.6);
    \draw[fill=black!8]  (9.75,0)  rectangle (10.85,0.6);
    \draw[fill=lenfill]  (10.85,0) rectangle (11.6,0.6);
    \draw (7.9,0) rectangle (11.6,0.6);
    \node at (8.575,0.3) {\textsf{C}};
    \node[font=\bfseries] at (9.5,0.3) {$\mathsf{1}$};
    \node at (10.3,0.3)  {$\mathsf{0}\cdots\mathsf{0}$};
    \node at (11.22,0.3) {$\langle \ell \rangle$};
    % pad bit and zero-fill annotations (staggered to avoid collision)
    \node[font=\scriptsize, align=center] (onebit) at (8.15,1.7)
      {single pad bit $\mathsf{1}$\\ \tiny(mandatory: exactly one, always
       present, marks end of message)};
    \draw[->, thin] (onebit.south) .. controls (8.9,1.15) .. (9.5,0.63);
    \node[font=\scriptsize, align=center] (zeros) at (10.9,1.15)
      {zero fill\\ \tiny($k$ zeros, $0\le k\le 511$; \emph{possibly none})};
    \draw[->, thin] (zeros.south west) -- (10.3,0.63);
    % ---- block names and sizes ----
    \node at (1.85,-0.33) {$m_1$};
    \node at (5.8,-0.33)  {$m_2$};
    \node at (9.4,-0.33) {$m_3$};
    \draw[decorate, decoration={brace, mirror, amplitude=4pt}]
      (0,-0.62) -- (3.7,-0.62)
      node[midway, below=5pt, font=\scriptsize] {512 bits};
    \draw[decorate, decoration={brace, mirror, amplitude=3pt}]
      (10.85,-0.08) -- (11.6,-0.08)
      node[midway, below=4pt, font=\scriptsize, align=center]
      {$64$ bits:\\bit length $\ell$};
  \end{tikzpicture}
  \caption{Padding and blocking, for a message needing $k=3$ blocks. The
  linear message (blue; portions \textsf{A}, \textsf{B}, \textsf{C})
  is consumed left to right by $512$-bit blocks. The final block is
  completed by a single $\mathsf{1}$ bit (green), a run of $\mathsf{0}$s,
  and the $64$-bit binary encoding of the message length $\ell$
  (orange)---the Merkle--Damg{\aa}rd strengthening. The rule is exact and
  admits no options: to a message of $\ell$ bits one \emph{always} appends
  one $\mathsf{1}$, then the fewest $\mathsf{0}$s (possibly none) so that
  the total reaches a multiple of $512$ once the $64$-bit length is
  affixed---i.e.\ $k$ is the least $\ge 0$ solution of
  $\ell + 1 + k + 64 \equiv 0 \pmod{512}$. The $\mathsf{1}$ bit and the
  length field are never omitted; only the zero-count $k$ depends on
  $\ell$, and if the tail \textsf{C} leaves no room for the $65$ closing
  bits, $k$ simply runs on into an extra all-padding block.}
  \label{fig:padding}
\end{figure}

\paragraph{Chaining.} All real work is delegated to one
\emph{compression function}
\[
  f : \State \times \bitstrings^{512} \longrightarrow \State,
  \qquad \State = (\Ztwo)^8 ,
\]
which eats a $512$-bit block and updates a point of the finite space
$\State$ (about which everything in \cref{sec:statespace}). Fix a public
starting point $h_0 = \mathrm{IV} \in \State$ and set
\[
  h_i = f(h_{i-1}, m_i), \qquad \SHA(m) = h_k ,
\]
reading the final state out as $256$ bits (\cref{fig:md}).

\begin{figure}[ht]
  \centering
  \begin{tikzpicture}[
      node distance=9mm,
      blk/.style={draw, minimum width=13mm, minimum height=7mm, font=\small},
      f/.style={draw, fill=projfill, minimum size=8mm, font=\small},
      >={Stealth[length=2mm]}]
    \node[f] (f1) {$f$};
    \node[f, right=16mm of f1] (f2) {$f$};
    \node[right=10mm of f2, font=\small] (dots) {$\cdots$};
    \node[f, right=10mm of dots] (f3) {$f$};
    \node[blk, above=7mm of f1] (m1) {$m_1$};
    \node[blk, above=7mm of f2] (m2) {$m_2$};
    \node[blk, above=7mm of f3, fill=black!8] (m3) {$m_k$};
    \node[left=9mm of f1, font=\small] (iv) {$\mathrm{IV}$};
    \node[right=9mm of f3, font=\small] (out) {$\SHA(m)$};
    \draw[->] (iv) -- (f1);
    \draw[->] (f1) -- node[below, font=\scriptsize]{$h_1$} (f2);
    \draw[->] (f2) -- (dots);
    \draw[->] (dots) -- node[below, font=\scriptsize]{$h_{k-1}$} (f3);
    \draw[->] (f3) -- (out);
    \draw[->] (m1) -- (f1);
    \draw[->] (m2) -- (f2);
    \draw[->] (m3) -- (f3);
    \node[above=1mm of m3, font=\scriptsize, align=center]
      {last block contains\\the bit length $\ell$};
  \end{tikzpicture}
  \caption{The Merkle--Damg{\aa}rd construction: an arbitrary-length
  message is folded through one compression function
  $f : \State \times \bitstrings^{512} \to \State$ on the finite state
  space $\State = (\Ztwo)^8$.}
  \label{fig:md}
\end{figure}

This is not an ad hoc recipe; it is a construction with a theorem, proved
independently in 1989 by Damg{\aa}rd---stated and proved as Theorem~3.1
of~\citet[pp.~419--420]{Damgard1989}---and by Merkle, under the name
``meta-method,'' in~\citet{MerkleOWHF1989}. Both original papers are
short, readable, and archived alongside this note; the theorem is exactly
why the design has been reused by MD5, SHA-1, and SHA-2 alike:

\begin{theorem}[Merkle, Damg{\aa}rd]\label{thm:md}
Any collision of the iterated hash yields, constructively, a collision of
the compression function: from $m \ne m'$ with $\SHA(m) = \SHA(m')$ one
can extract $(h, b) \ne (h', b')$ with $f(h,b) = f(h',b')$.
\end{theorem}

\begin{proof}[Proof idea]
Walk the two chains backward from the equal outputs. If the messages have
different lengths, the final blocks already differ (they encode the
lengths---this is where strengthening is used) yet map to the same output:
a collision in $f$. If the lengths agree, compare positions from the end;
at the last position where the inputs to $f$ differ, the outputs agree,
again a collision in $f$.
\end{proof}

The moral: \emph{the arbitrary-length-input problem reduces, with proof,
to the study of one map on one finite space}. (The domain is itself
finite---every allowed input is shorter than $2^{64}$ bits---but its
inputs run through unboundedly many \emph{lengths}, and it is that
length-variability, not any actual infinity, that the construction tames.)
The reduction is airtight only for
collision resistance---the construction leaks other structure, such as the
famous length-extension property (from the digest $H(m)$ alone, without
knowing $m$, one can compute $H(m \Vert \mathrm{pad} \Vert m')$; defined
and dissected in \cref{sec:models}), whose repair spawned its own theory
(indifferentiability) and whose avoidance shaped
the post-2000 designs, including the sponge construction underlying
SHA-3~\citep{BDPV2007}. But it means the rest of this paper may honestly
focus on the finite object $f$.

\subsection{The state space and its two rival geometries}
\label{sec:statespace}

The state space is
\[
  \State \;=\; (\Ztwo)^{8},
  \qquad |\State| = 2^{256} :
\]
an $8$-dimensional module-like array of $32$-bit words---the
``$8{\times}32$'' geometry in which the entire life of \SHA{} unfolds. As
a plain set, $\State$ is also $\Ftwo^{256}$, a $256$-dimensional vector
space over the two-element field. Everything hinges on the fact that
\emph{the same underlying set carries two incompatible abelian group
structures}, both used incessantly:
\begin{itemize}
  \item bitwise addition $\shaxor$ (exclusive or): the $\Ftwo$-vector
  space structure of $\Ftwo^{256}$;
  \item wordwise addition $\shaplus$ modulo $2^{32}$: the product group
  structure of $(\Ztwo)^8$, where carries propagate within each word.
\end{itemize}
These two structures share the zero element and almost nothing else. A map
that is linear for $\shaxor$ (a matrix over $\Ftwo$) is, generically,
maximally nonlinear from the point of view of $\shaplus$, and vice versa;
the carry chains of integer addition are precisely the terms that
$\Ftwo$-linear algebra cannot see.

It helps to picture the additive ($\shaplus$) structure literally, as a
clock (\cref{fig:odometer}). Read the state as an odometer counting
modulo $2^{256}$: every modular addition advances the needle by some
amount, and the crucial event is whether it \emph{wraps} past $0$---an
overflow, a carry out of the top, thrown away by the reduction. That one
bit of ``did it go around?'' is exactly what $\shaxor$-algebra is blind to
and what the $2$-adic metric, next, is built to record.

\begin{figure}[ht]
  \centering
  \begin{tikzpicture}[font=\small, >={Stealth[length=2mm]}]
    \def\R{1.75}
    \draw[thick] (0,0) circle (\R);
    \fill (90:\R) circle (1.7pt);
    \node[above=1pt] at (90:\R) {\scriptsize $0 \equiv 2^{256}$};
    \fill (-90:\R) circle (1.3pt);
    \node[below=1pt] at (-90:\R) {\scriptsize $2^{255}$};
    \node[align=center, font=\scriptsize] at (0,0)
      {state as a\\ counter\\ $\bmod\; 2^{256}$};
    % arc 1: no wrap (does not cross the top marker at 90)
    \fill[projframe] (25:\R) circle (1.5pt);
    \draw[->, thick, projframe] (25:{\R+0.3}) arc (25:-120:{\R+0.3});
    \node[projframe, font=\scriptsize, align=center] at (-52:{\R+1.25})
      {$x \shaplus c_1$:\\ stays short of $2^{256}$\\ \emph{no carry}};
    % arc 2: wraps past the top marker (crosses 90)
    \fill[citewine] (145:\R) circle (1.5pt);
    \draw[->, thick, citewine] (145:{\R+0.3}) arc (145:35:{\R+0.3});
    \node[citewine, font=\scriptsize, align=center] at (130:{\R+1.5})
      {$x \shaplus c_2$: sweeps\\ past $0$ --- \emph{carry out}\\ (the wrap)};
  \end{tikzpicture}
  \caption{Wordwise addition $\shaplus$ as motion on a clock. As an
  odometer counting mod $2^{256}$, the state advances by each addition;
  it \emph{wraps} past $0$ exactly when the sum overflows---a carry out of
  the top bit, discarded by the reduction. Whether a given addition wraps
  is invisible to $\Ftwo$-linear ($\shaxor$) algebra and is precisely the
  information the $2$-adic metric records. (The true state is not one
  circle of $2^{256}$ positions but eight \emph{coupled} clocks of
  $2^{32}$ each, $(\Ztwo)^8$: the single odometer is the intuition, the
  product of eight is the reality.) A round is a great many such
  additions, and iterating \SHA{} (\cref{sec:iterated}) sends the counter
  wandering---sometimes around the circle, sometimes not.}
  \label{fig:odometer}
\end{figure}

There is a third, more analytic way to say this, and it is the one we
expect resonates with a $p$-adic sensibility. Identify $\Ztwo$ with the
quotient $\mathbb{Z}_2 / 2^{32}\mathbb{Z}_2$ of the $2$-adic integers.
Then the basic machine operations sort themselves by how they treat the
$2$-adic metric (in which two integers are close when they agree in many
low-order bits):
\begin{center}
  \begin{tabular}{@{}lccc@{}}
    \toprule
    operation & $\Ftwo$-linear? & $2$-adically $1$-Lipschitz? & where it is ``tame'' \\
    \midrule
    $x \shaxor c$, bitwise $\wedge,\vee,\lnot$ & yes ($\shaxor$) / no & yes & both lenses see it \\
    $x \shaplus c$, \ $x \shaplus y$          & no  & yes & only the $2$-adic lens \\
    rotation $x \lll r$, shift $x \gg s$       & yes & no  & only the $\Ftwo$ lens \\
    \bottomrule
  \end{tabular}
\end{center}
Bitwise logic and modular addition are \emph{$1$-Lipschitz} ($2$-adically
continuous in the strongest sense: output bit $i$ depends only on input
bits $0,\dots,i$, the triangular or ``causal'' maps); compositions of
these are the \emph{T-functions} of Klimov and
Shamir~\citep{KlimovShamir2003}, and there is a genuine ergodic theory,
due to Anashin, characterizing when such maps act
measure-preservingly or ergodically on $\mathbb{Z}_2$~\citep{Anashin2006}.
Rotations, by contrast, are $\Ftwo$-linear but violently discontinuous
$2$-adically: they let high-order bits influence low-order ones.

The $2$-adic lens is not an ad hoc bookkeeping device; it opens onto a
genuine subject. Nonarchimedean dynamics---the iteration theory of
self-maps of $p$-adic disks, opened up by \citet{Lubin1994}---asks
exactly this paper's kind of question one level of structure up: when
must the iteration of a $p$-adic map secretly respect algebraic
structure? Lubin's paper records a suspicion of exactly the Josephus
flavor---that a noninvertible series commuting with an invertible
nontorsion one must have a formal group hidden behind it---and its main
result makes a sharp instance of that suspicion into a theorem
(\citealp{Lubin1994}, Main Theorem~6.3):

\begin{theorem}[Lubin]\label{thm:lubin}
Let $o$ be the ring of integers in a finite extension of $\mathbb{Q}_p$,
with maximal ideal $\mathfrak m$, and let $u, f \in o[[x]]$ be power
series fixing $0$, with $u$ \emph{invertible} (its linear coefficient
$u'(0)$ a unit) and $f$ \emph{noninvertible} ($f'(0)\in\mathfrak m$) of
finite Weierstrass degree $\delta$---the number of zeros $f$ has in the
open disk, counted with multiplicity. If $u$ and $f$ commute,
$u\circ f = f\circ u$, then \emph{either} some iterate of $u$ is the
identity series ($u$ is a root of unity in the group of invertible
series) \emph{or} $\delta = p^{d}$ for some $d\ge 1$.
\end{theorem}

The content is rigidity. The only Weierstrass degrees a noninvertible
map may carry while sharing its orbits with a genuine (nontorsion)
invertible symmetry are the powers of $p$---and those are precisely the
degrees of the endomorphisms of a one-dimensional formal group. So bare
commutation, a purely dynamical hypothesis with no algebra in it, already
forces the numerical signature of a formal-group skeleton; whether the
skeleton can ever be truly absent is Lubin's own still-open question.

The phenomenon is visible in miniature over our own ring $\mathbb{Z}_2$
($p=2$). The noninvertible $f(x) = 4x + x^{2}$---linear coefficient
$4\in 2\mathbb{Z}_2$, Weierstrass degree $2 = 2^{1}$, a $p$-power on the
nose---commutes with the invertible $u(x) = 9x + 6x^{2} + x^{3}$ (unit
linear coefficient $9$), as one checks by expanding
$u\circ f = f\circ u$ directly. The companion Corollary~3.2.1 of the same
paper is just as concrete: every zero of $f'$ inside the disk must be a
zero of some iterate of $f$. Here $f'(x) = 2x + 4$ vanishes at $x=-2$,
and indeed $f(f(-2)) = f(-4) = 0$. Two lines of arithmetic thus exhibit a
symmetry ($u$), a forced degree ($2^{1}$), and a forced coincidence among
the iterates---structure precipitating out of commutation alone, the
Josephus moral in $p$-adic dress.

How near does this reach \SHA{}? Adjacent, honestly, not overhead.
\Cref{thm:lubin} governs \emph{analytic} series on the $p$-adic disk,
whereas \SHA{}'s ARX pieces are merely $1$-Lipschitz T-functions---the
coarse, non-analytic relatives tabulated above---so the theorem does not
bind the function on the nose, and Anashin's ergodic theory is what
survives at that lower regularity~\citep{Anashin2006, KlimovShamir2003}.
What \emph{does} transfer is the question and its shape: does any
nontrivial map commute with a (reduced, word-narrowed) \SHA{} component,
and if one did, would it drag in hidden algebraic structure the way it
must for analytic series? At toy scale that is a finite, machine-checkable
hunt---search for a commuting symmetry; find none, and the absence is
itself a measurement of structurelessness. In this language \SHA{}'s
designers are in the business of manufacturing $1$-Lipschitz dynamics
that commute with nothing and hide no formal group, and \cref{thm:lubin}
is the evidence that, one regularity class up, such structurelessness is
genuinely hard to achieve by accident.

\SHA{}'s round function is an \emph{ARX} design---Addition, Rotation,
Xor---and the table explains the acronym's cunning: each primitive is
transparent in one geometry and wild in the other, and the design
alternates them so that the composition is wild in both. An analyst who
adopts $\Ftwo$-linear algebra kills the X's and R's but is left helpless
against the carries; an analyst who adopts the $2$-adic/T-function lens
domesticates additions and logic but is destroyed by the rotations. This
is not an accident of \SHA{}; it is the design principle of the whole ARX
family. What \emph{is} specific to \SHA{} is the exact schedule of
alternation, and \cref{sec:falling} showed how much such specifics matter.

\begin{project}{Single operations under the $2$-adic lens}{elementary
  number theory; curiosity about $p$-adic numbers}{3--4 weeks}
  Anashin's theory gives clean criteria for when a T-function like
  $x \mapsto x \shaplus c$ or $x \mapsto x \shaxor c \shaplus (x \wedge d)$
  is a bijection of $\Ztwo[k]$ for all $k$, and when its iteration visits
  every residue (ergodicity)~\citep{Anashin2006, KlimovShamir2003}.
  Verify the criteria experimentally on $8$- and $16$-bit words, visualize
  orbits, then break the hypotheses with a single rotation and measure
  what dies. A perfect first encounter with $p$-adic ideas producing
  visible, computable phenomena.
\end{project}

\subsection{The compression function, structurally}
\label{sec:roundshape}

We describe $f$ at the level of shape; the full specification fits on two
pages of \citet{FIPS180-4} (with an accessible pseudocode rendering
in~\citealp{WikipediaSHA2}) and is worth reading only after the shape is
clear. The state is eight words $(a,b,c,d,e,f,g,h)$. The $512$-bit message
block is first \emph{expanded}: its sixteen words $W_0, \dots, W_{15}$
are stretched to sixty-four by the recurrence
\[
  W_t \;=\; \sigma_1(W_{t-2}) \shaplus W_{t-7} \shaplus
            \sigma_0(W_{t-15}) \shaplus W_{t-16},
  \qquad 16 \le t < 64,
\]
where $\sigma_0, \sigma_1$ are fixed $\shaxor$-combinations of rotations
and shifts of a single word ($\Ftwo$-linear), while the additions are
modular (2-adic). Then sixty-four rounds are applied; round $t$ computes
\[
  T_1 = h \shaplus \Sigma_1(e) \shaplus \mathrm{Ch}(e,f,g)
          \shaplus K_t \shaplus W_t,
  \qquad
  T_2 = \Sigma_0(a) \shaplus \mathrm{Maj}(a,b,c),
\]
and shifts the state pipeline:
$(a, \dots, h) \leftarrow
 (T_1 \shaplus T_2,\; a,\; b,\; c,\; d \shaplus T_1,\; e,\; f,\; g)$.
Here $\Sigma_0, \Sigma_1$ are again xors of three rotations, and
\[
  \mathrm{Ch}(x,y,z) = (x \wedge y) \shaxor (\lnot x \wedge z),
  \qquad
  \mathrm{Maj}(x,y,z) = (x \wedge y) \shaxor (x \wedge z) \shaxor (y \wedge z)
\]
are bitwise Boolean functions of degree two (``choose'' and ``majority'').
Finally---a crucial asymmetry called the Davies--Meyer feedforward---the
input state is added back: the block's net effect is
$h \mapsto h \shaplus \tilde{f}(h, m)$, which prevents running the rounds
backward even though each individual round is invertible. This
feedforward is not a detail but the hinge on which one-wayness turns, and
it is one of the two places where \SHA{}'s architecture carries a genuine
theorem; we return to it in \cref{sec:constructions}.

\Cref{fig:schedule} traces how a single block's data reaches every one of
the sixty-four rounds. The block supplies only the first sixteen schedule
words; the recurrence then \emph{stretches} those seeds forward, each new
word reaching back to four earlier ones, until all sixty-four exist---so a
bit set in the incoming block propagates, through the schedule alone,
into arbitrarily late rounds. Each round $t$ then draws exactly two
external inputs, its schedule word $W_t$ and its constant $K_t$, and folds
them into the running eight-word state; everything else in the round is
internal feedback.

\begin{figure}[ht]
  \centering
  \begin{tikzpicture}[font=\small, >={Stealth[length=1.7mm]},
      seed/.style={draw, fill=msgfill, minimum width=3.1mm,
                   minimum height=5mm, inner sep=0pt},
      der/.style={draw, fill=projfill, minimum width=3.1mm,
                  minimum height=5mm, inner sep=0pt},
      rnd/.style={draw, fill=projfill, minimum width=7mm,
                  minimum height=6mm, inner sep=0pt, font=\scriptsize},
      lbl/.style={font=\scriptsize}]
    % ---- the 512-bit block ----
    \draw[fill=msgfill] (0,5.0) rectangle (5.28,5.5);
    \node[lbl] at (2.64,5.25) {$512$-bit message block $m$};
    % ---- 16 seed words W0..W15 ----
    \foreach \i in {0,...,15} {
      \node[seed] (w\i) at ({\i*0.33+0.165},4.2) {};
      \draw[->, gray] ({\i*0.33+0.165},4.98) -- (w\i.north);
    }
    \node[lbl, below=0pt of w0] {$W_0$};
    \node[lbl, below=0pt of w15] {$W_{15}$};
    \node[lbl, align=center] at (2.5,3.55) {\scriptsize\emph{16 seed words (from block)}};
    % ---- derived words W16 .. W63 ----
    \node[der, fill=lenfill] (w16) at (6.9,4.2) {};
    \node[lbl, below=0pt of w16] {$W_{16}$};
    \node[der] (w17) at (7.25,4.2) {};
    \node[lbl] at (7.75,4.2) {$\cdots$};
    \node[der] (w63) at (8.35,4.2) {};
    \node[lbl, below=0pt of w63] {$W_{63}$};
    \node[lbl, align=center] at (7.6,3.55) {\scriptsize\emph{48 derived}};
    % fan-in into W16 from W0,W1(sigma0),W9,W14(sigma1)
    \draw[->, projframe] (w0.north) ++(0,0.02) to[bend left=32] (w16.north west);
    \draw[->, projframe] (w1.north) to[bend left=30]
      node[above, lbl, pos=0.62] {$\sigma_0$} (w16.north);
    \draw[->, projframe] (w9.north) to[bend left=22] (w16.north);
    \draw[->, projframe] (w14.north) to[bend left=16]
      node[above, lbl, pos=0.5] {$\sigma_1$} (w16.north east);
    \node[lbl, align=center, anchor=south] at (3.2,5.95)
      {each derived word reaches back four:
       $W_t=\sigma_1(W_{t-2})\shaplus W_{t-7}\shaplus\sigma_0(W_{t-15})\shaplus W_{t-16}$};
    % ---- bridge arrow ----
    \node[lbl, anchor=east] at (9.9,3.35)
      {\emph{the $64$ words feed the $64$ rounds, one each}};
    \draw[->, thick] (5.0,3.55) -- (5.0,2.75);
    % ---- round pipeline ----
    \node[lbl, anchor=east] at (-0.2,1.9) {$\mathrm{IV}$};
    \draw[->, very thick, black!55] (-0.2,1.9) -- (0.28,1.9);
    \foreach \i/\lab/\x in {0/{$R_0$}/0.9, 1/{$R_1$}/2.35, 2/{$R_t$}/4.8,
                            3/{$R_{63}$}/7.4} {
      \node[rnd, minimum width=9mm] (r\i) at (\x,1.9) {\lab};
    }
    \node[lbl] at (3.575,1.9) {$\cdots$};
    \node[lbl] at (6.1,1.9) {$\cdots$};
    \draw[->, very thick, black!55] (r0.east) -- (r1.west);
    \draw[->, very thick, black!55] (r1.east) -- (3.2,1.9);
    \draw[->, very thick, black!55] (3.95,1.9) -- (r2.west);
    \draw[->, very thick, black!55] (r2.east) -- (5.75,1.9);
    \draw[->, very thick, black!55] (6.45,1.9) -- (r3.west);
    \draw[->, very thick, black!55] (r3.east) -- (8.15,1.9)
      node[lbl, anchor=west, black] {$h_{\mathrm{out}}$};
    \node[lbl, black!55, anchor=north] at (3.7,1.42)
      {the eight-word state $(a,\dots,h)$, threaded left to right};
    % W (wine) and K (navy) inputs into each shown round
    \foreach \r/\wl/\kl in {r0/{$W_0$}/{$K_0$}, r1/{$W_1$}/{$K_1$},
                            r2/{$W_t$}/{$K_t$}, r3/{$W_{63}$}/{$K_{63}$}} {
      \draw[->, citewine] ($(\r.north)+(-0.28,0.6)$) -- ($(\r.north)+(-0.28,0.02)$);
      \node[lbl, citewine, anchor=south] at ($(\r.north)+(-0.28,0.58)$) {\wl};
      \draw[->, linknavy] ($(\r.north)+(0.28,0.6)$) -- ($(\r.north)+(0.28,0.02)$);
      \node[lbl, linknavy, anchor=south] at ($(\r.north)+(0.28,0.58)$) {\kl};
    }
  \end{tikzpicture}
  \caption{How one block's data reaches every round. The block fills only
  the sixteen seed words $W_0,\dots,W_{15}$ (blue); the message-schedule
  recurrence then stretches them into $W_{16},\dots,W_{63}$ (each derived
  word, shown for $W_{16}$, built from four earlier ones via the
  $\Ftwo$-linear maps $\sigma_0,\sigma_1$ and modular addition). Each of
  the $64$ rounds $R_t$ consumes exactly two external inputs---its
  schedule word $W_t$ (wine) and its constant $K_t$ (navy)---and threads
  the eight-word state through from the incoming chaining value to
  $h_{\mathrm{out}}$. Arrows show influence: a bit of the block, via the
  schedule's reach-back, can steer arbitrarily late rounds.}
  \label{fig:schedule}
\end{figure}

The reach of that schedule is worth making quantitative, because it is a
clean computed pattern rather than a hand-drawn cartoon. Trace, for each
word $W_t$, the set of \emph{seed} words $W_0,\dots,W_{15}$---the only
ones carrying message data---that it transitively depends on, following
the recurrence's four back-pointers ($t-2,\,t-7,\,t-15,\,t-16$) down to
the data layer. The dependency graph
(\cref{fig:schedulegraph}) then answers ``how fast does a block's data
spread through the schedule?'' exactly: the count of seeds reached climbs
$4,4,7,7,10,10,13,14,15,15,\dots$ and \emph{saturates at $W_{26}$}---from
that word on, every schedule word depends on all sixteen data words. Ten
steps of the recurrence suffice to fully entangle the message, before a
single one of the sixty-four rounds has done its diffusion work. The
four coloured arc families in the figure are the four terms of the
recurrence, and their regular weave is the schedule's linear-feedback
skeleton laid bare; whether that visible regularity is a foothold or a
red herring is one more question small experiments can chew on.

\begin{figure}[ht]
  \centering
  \resizebox{\linewidth}{!}{% AUTO-GENERATED by tools/wschedule_graph.py -- do not edit by hand.
\begin{tikzpicture}[font=\scriptsize, >={Stealth[length=1.4mm]}]
  \def\dx{0.205}
  \draw[citewine, line width=0.2pt, opacity=0.4] (2.870,0) to[out=58,in=122] (3.280,0);
  \draw[projframe, line width=0.2pt, opacity=0.4] (1.845,0) to[out=58,in=122] (3.280,0);
  \draw[linknavy, line width=0.2pt, opacity=0.4] (0.205,0) to[out=58,in=122] (3.280,0);
  \draw[orange!80!black, line width=0.2pt, opacity=0.4] (0.000,0) to[out=58,in=122] (3.280,0);
  \draw[citewine, line width=0.2pt, opacity=0.4] (3.075,0) to[out=58,in=122] (3.485,0);
  \draw[projframe, line width=0.2pt, opacity=0.4] (2.050,0) to[out=58,in=122] (3.485,0);
  \draw[linknavy, line width=0.2pt, opacity=0.4] (0.410,0) to[out=58,in=122] (3.485,0);
  \draw[orange!80!black, line width=0.2pt, opacity=0.4] (0.205,0) to[out=58,in=122] (3.485,0);
  \draw[citewine, line width=0.2pt, opacity=0.4] (3.280,0) to[out=58,in=122] (3.690,0);
  \draw[projframe, line width=0.2pt, opacity=0.4] (2.255,0) to[out=58,in=122] (3.690,0);
  \draw[linknavy, line width=0.2pt, opacity=0.4] (0.615,0) to[out=58,in=122] (3.690,0);
  \draw[orange!80!black, line width=0.2pt, opacity=0.4] (0.410,0) to[out=58,in=122] (3.690,0);
  \draw[citewine, line width=0.2pt, opacity=0.4] (3.485,0) to[out=58,in=122] (3.895,0);
  \draw[projframe, line width=0.2pt, opacity=0.4] (2.460,0) to[out=58,in=122] (3.895,0);
  \draw[linknavy, line width=0.2pt, opacity=0.4] (0.820,0) to[out=58,in=122] (3.895,0);
  \draw[orange!80!black, line width=0.2pt, opacity=0.4] (0.615,0) to[out=58,in=122] (3.895,0);
  \draw[citewine, line width=0.2pt, opacity=0.4] (3.690,0) to[out=58,in=122] (4.100,0);
  \draw[projframe, line width=0.2pt, opacity=0.4] (2.665,0) to[out=58,in=122] (4.100,0);
  \draw[linknavy, line width=0.2pt, opacity=0.4] (1.025,0) to[out=58,in=122] (4.100,0);
  \draw[orange!80!black, line width=0.2pt, opacity=0.4] (0.820,0) to[out=58,in=122] (4.100,0);
  \draw[citewine, line width=0.2pt, opacity=0.4] (3.895,0) to[out=58,in=122] (4.305,0);
  \draw[projframe, line width=0.2pt, opacity=0.4] (2.870,0) to[out=58,in=122] (4.305,0);
  \draw[linknavy, line width=0.2pt, opacity=0.4] (1.230,0) to[out=58,in=122] (4.305,0);
  \draw[orange!80!black, line width=0.2pt, opacity=0.4] (1.025,0) to[out=58,in=122] (4.305,0);
  \draw[citewine, line width=0.2pt, opacity=0.4] (4.100,0) to[out=58,in=122] (4.510,0);
  \draw[projframe, line width=0.2pt, opacity=0.4] (3.075,0) to[out=58,in=122] (4.510,0);
  \draw[linknavy, line width=0.2pt, opacity=0.4] (1.435,0) to[out=58,in=122] (4.510,0);
  \draw[orange!80!black, line width=0.2pt, opacity=0.4] (1.230,0) to[out=58,in=122] (4.510,0);
  \draw[citewine, line width=0.2pt, opacity=0.4] (4.305,0) to[out=58,in=122] (4.715,0);
  \draw[projframe, line width=0.2pt, opacity=0.4] (3.280,0) to[out=58,in=122] (4.715,0);
  \draw[linknavy, line width=0.2pt, opacity=0.4] (1.640,0) to[out=58,in=122] (4.715,0);
  \draw[orange!80!black, line width=0.2pt, opacity=0.4] (1.435,0) to[out=58,in=122] (4.715,0);
  \draw[citewine, line width=0.2pt, opacity=0.4] (4.510,0) to[out=58,in=122] (4.920,0);
  \draw[projframe, line width=0.2pt, opacity=0.4] (3.485,0) to[out=58,in=122] (4.920,0);
  \draw[linknavy, line width=0.2pt, opacity=0.4] (1.845,0) to[out=58,in=122] (4.920,0);
  \draw[orange!80!black, line width=0.2pt, opacity=0.4] (1.640,0) to[out=58,in=122] (4.920,0);
  \draw[citewine, line width=0.2pt, opacity=0.4] (4.715,0) to[out=58,in=122] (5.125,0);
  \draw[projframe, line width=0.2pt, opacity=0.4] (3.690,0) to[out=58,in=122] (5.125,0);
  \draw[linknavy, line width=0.2pt, opacity=0.4] (2.050,0) to[out=58,in=122] (5.125,0);
  \draw[orange!80!black, line width=0.2pt, opacity=0.4] (1.845,0) to[out=58,in=122] (5.125,0);
  \draw[citewine, line width=0.2pt, opacity=0.4] (4.920,0) to[out=58,in=122] (5.330,0);
  \draw[projframe, line width=0.2pt, opacity=0.4] (3.895,0) to[out=58,in=122] (5.330,0);
  \draw[linknavy, line width=0.2pt, opacity=0.4] (2.255,0) to[out=58,in=122] (5.330,0);
  \draw[orange!80!black, line width=0.2pt, opacity=0.4] (2.050,0) to[out=58,in=122] (5.330,0);
  \draw[citewine, line width=0.2pt, opacity=0.4] (5.125,0) to[out=58,in=122] (5.535,0);
  \draw[projframe, line width=0.2pt, opacity=0.4] (4.100,0) to[out=58,in=122] (5.535,0);
  \draw[linknavy, line width=0.2pt, opacity=0.4] (2.460,0) to[out=58,in=122] (5.535,0);
  \draw[orange!80!black, line width=0.2pt, opacity=0.4] (2.255,0) to[out=58,in=122] (5.535,0);
  \draw[citewine, line width=0.2pt, opacity=0.4] (5.330,0) to[out=58,in=122] (5.740,0);
  \draw[projframe, line width=0.2pt, opacity=0.4] (4.305,0) to[out=58,in=122] (5.740,0);
  \draw[linknavy, line width=0.2pt, opacity=0.4] (2.665,0) to[out=58,in=122] (5.740,0);
  \draw[orange!80!black, line width=0.2pt, opacity=0.4] (2.460,0) to[out=58,in=122] (5.740,0);
  \draw[citewine, line width=0.2pt, opacity=0.4] (5.535,0) to[out=58,in=122] (5.945,0);
  \draw[projframe, line width=0.2pt, opacity=0.4] (4.510,0) to[out=58,in=122] (5.945,0);
  \draw[linknavy, line width=0.2pt, opacity=0.4] (2.870,0) to[out=58,in=122] (5.945,0);
  \draw[orange!80!black, line width=0.2pt, opacity=0.4] (2.665,0) to[out=58,in=122] (5.945,0);
  \draw[citewine, line width=0.2pt, opacity=0.4] (5.740,0) to[out=58,in=122] (6.150,0);
  \draw[projframe, line width=0.2pt, opacity=0.4] (4.715,0) to[out=58,in=122] (6.150,0);
  \draw[linknavy, line width=0.2pt, opacity=0.4] (3.075,0) to[out=58,in=122] (6.150,0);
  \draw[orange!80!black, line width=0.2pt, opacity=0.4] (2.870,0) to[out=58,in=122] (6.150,0);
  \draw[citewine, line width=0.2pt, opacity=0.4] (5.945,0) to[out=58,in=122] (6.355,0);
  \draw[projframe, line width=0.2pt, opacity=0.4] (4.920,0) to[out=58,in=122] (6.355,0);
  \draw[linknavy, line width=0.2pt, opacity=0.4] (3.280,0) to[out=58,in=122] (6.355,0);
  \draw[orange!80!black, line width=0.2pt, opacity=0.4] (3.075,0) to[out=58,in=122] (6.355,0);
  \draw[citewine, line width=0.2pt, opacity=0.4] (6.150,0) to[out=58,in=122] (6.560,0);
  \draw[projframe, line width=0.2pt, opacity=0.4] (5.125,0) to[out=58,in=122] (6.560,0);
  \draw[linknavy, line width=0.2pt, opacity=0.4] (3.485,0) to[out=58,in=122] (6.560,0);
  \draw[orange!80!black, line width=0.2pt, opacity=0.4] (3.280,0) to[out=58,in=122] (6.560,0);
  \draw[citewine, line width=0.2pt, opacity=0.4] (6.355,0) to[out=58,in=122] (6.765,0);
  \draw[projframe, line width=0.2pt, opacity=0.4] (5.330,0) to[out=58,in=122] (6.765,0);
  \draw[linknavy, line width=0.2pt, opacity=0.4] (3.690,0) to[out=58,in=122] (6.765,0);
  \draw[orange!80!black, line width=0.2pt, opacity=0.4] (3.485,0) to[out=58,in=122] (6.765,0);
  \draw[citewine, line width=0.2pt, opacity=0.4] (6.560,0) to[out=58,in=122] (6.970,0);
  \draw[projframe, line width=0.2pt, opacity=0.4] (5.535,0) to[out=58,in=122] (6.970,0);
  \draw[linknavy, line width=0.2pt, opacity=0.4] (3.895,0) to[out=58,in=122] (6.970,0);
  \draw[orange!80!black, line width=0.2pt, opacity=0.4] (3.690,0) to[out=58,in=122] (6.970,0);
  \draw[citewine, line width=0.2pt, opacity=0.4] (6.765,0) to[out=58,in=122] (7.175,0);
  \draw[projframe, line width=0.2pt, opacity=0.4] (5.740,0) to[out=58,in=122] (7.175,0);
  \draw[linknavy, line width=0.2pt, opacity=0.4] (4.100,0) to[out=58,in=122] (7.175,0);
  \draw[orange!80!black, line width=0.2pt, opacity=0.4] (3.895,0) to[out=58,in=122] (7.175,0);
  \draw[citewine, line width=0.2pt, opacity=0.4] (6.970,0) to[out=58,in=122] (7.380,0);
  \draw[projframe, line width=0.2pt, opacity=0.4] (5.945,0) to[out=58,in=122] (7.380,0);
  \draw[linknavy, line width=0.2pt, opacity=0.4] (4.305,0) to[out=58,in=122] (7.380,0);
  \draw[orange!80!black, line width=0.2pt, opacity=0.4] (4.100,0) to[out=58,in=122] (7.380,0);
  \draw[citewine, line width=0.2pt, opacity=0.4] (7.175,0) to[out=58,in=122] (7.585,0);
  \draw[projframe, line width=0.2pt, opacity=0.4] (6.150,0) to[out=58,in=122] (7.585,0);
  \draw[linknavy, line width=0.2pt, opacity=0.4] (4.510,0) to[out=58,in=122] (7.585,0);
  \draw[orange!80!black, line width=0.2pt, opacity=0.4] (4.305,0) to[out=58,in=122] (7.585,0);
  \draw[citewine, line width=0.2pt, opacity=0.4] (7.380,0) to[out=58,in=122] (7.790,0);
  \draw[projframe, line width=0.2pt, opacity=0.4] (6.355,0) to[out=58,in=122] (7.790,0);
  \draw[linknavy, line width=0.2pt, opacity=0.4] (4.715,0) to[out=58,in=122] (7.790,0);
  \draw[orange!80!black, line width=0.2pt, opacity=0.4] (4.510,0) to[out=58,in=122] (7.790,0);
  \draw[citewine, line width=0.2pt, opacity=0.4] (7.585,0) to[out=58,in=122] (7.995,0);
  \draw[projframe, line width=0.2pt, opacity=0.4] (6.560,0) to[out=58,in=122] (7.995,0);
  \draw[linknavy, line width=0.2pt, opacity=0.4] (4.920,0) to[out=58,in=122] (7.995,0);
  \draw[orange!80!black, line width=0.2pt, opacity=0.4] (4.715,0) to[out=58,in=122] (7.995,0);
  \draw[citewine, line width=0.2pt, opacity=0.4] (7.790,0) to[out=58,in=122] (8.200,0);
  \draw[projframe, line width=0.2pt, opacity=0.4] (6.765,0) to[out=58,in=122] (8.200,0);
  \draw[linknavy, line width=0.2pt, opacity=0.4] (5.125,0) to[out=58,in=122] (8.200,0);
  \draw[orange!80!black, line width=0.2pt, opacity=0.4] (4.920,0) to[out=58,in=122] (8.200,0);
  \draw[citewine, line width=0.2pt, opacity=0.4] (7.995,0) to[out=58,in=122] (8.405,0);
  \draw[projframe, line width=0.2pt, opacity=0.4] (6.970,0) to[out=58,in=122] (8.405,0);
  \draw[linknavy, line width=0.2pt, opacity=0.4] (5.330,0) to[out=58,in=122] (8.405,0);
  \draw[orange!80!black, line width=0.2pt, opacity=0.4] (5.125,0) to[out=58,in=122] (8.405,0);
  \draw[citewine, line width=0.2pt, opacity=0.4] (8.200,0) to[out=58,in=122] (8.610,0);
  \draw[projframe, line width=0.2pt, opacity=0.4] (7.175,0) to[out=58,in=122] (8.610,0);
  \draw[linknavy, line width=0.2pt, opacity=0.4] (5.535,0) to[out=58,in=122] (8.610,0);
  \draw[orange!80!black, line width=0.2pt, opacity=0.4] (5.330,0) to[out=58,in=122] (8.610,0);
  \draw[citewine, line width=0.2pt, opacity=0.4] (8.405,0) to[out=58,in=122] (8.815,0);
  \draw[projframe, line width=0.2pt, opacity=0.4] (7.380,0) to[out=58,in=122] (8.815,0);
  \draw[linknavy, line width=0.2pt, opacity=0.4] (5.740,0) to[out=58,in=122] (8.815,0);
  \draw[orange!80!black, line width=0.2pt, opacity=0.4] (5.535,0) to[out=58,in=122] (8.815,0);
  \draw[citewine, line width=0.2pt, opacity=0.4] (8.610,0) to[out=58,in=122] (9.020,0);
  \draw[projframe, line width=0.2pt, opacity=0.4] (7.585,0) to[out=58,in=122] (9.020,0);
  \draw[linknavy, line width=0.2pt, opacity=0.4] (5.945,0) to[out=58,in=122] (9.020,0);
  \draw[orange!80!black, line width=0.2pt, opacity=0.4] (5.740,0) to[out=58,in=122] (9.020,0);
  \draw[citewine, line width=0.2pt, opacity=0.4] (8.815,0) to[out=58,in=122] (9.225,0);
  \draw[projframe, line width=0.2pt, opacity=0.4] (7.790,0) to[out=58,in=122] (9.225,0);
  \draw[linknavy, line width=0.2pt, opacity=0.4] (6.150,0) to[out=58,in=122] (9.225,0);
  \draw[orange!80!black, line width=0.2pt, opacity=0.4] (5.945,0) to[out=58,in=122] (9.225,0);
  \draw[citewine, line width=0.2pt, opacity=0.4] (9.020,0) to[out=58,in=122] (9.430,0);
  \draw[projframe, line width=0.2pt, opacity=0.4] (7.995,0) to[out=58,in=122] (9.430,0);
  \draw[linknavy, line width=0.2pt, opacity=0.4] (6.355,0) to[out=58,in=122] (9.430,0);
  \draw[orange!80!black, line width=0.2pt, opacity=0.4] (6.150,0) to[out=58,in=122] (9.430,0);
  \draw[citewine, line width=0.2pt, opacity=0.4] (9.225,0) to[out=58,in=122] (9.635,0);
  \draw[projframe, line width=0.2pt, opacity=0.4] (8.200,0) to[out=58,in=122] (9.635,0);
  \draw[linknavy, line width=0.2pt, opacity=0.4] (6.560,0) to[out=58,in=122] (9.635,0);
  \draw[orange!80!black, line width=0.2pt, opacity=0.4] (6.355,0) to[out=58,in=122] (9.635,0);
  \draw[citewine, line width=0.2pt, opacity=0.4] (9.430,0) to[out=58,in=122] (9.840,0);
  \draw[projframe, line width=0.2pt, opacity=0.4] (8.405,0) to[out=58,in=122] (9.840,0);
  \draw[linknavy, line width=0.2pt, opacity=0.4] (6.765,0) to[out=58,in=122] (9.840,0);
  \draw[orange!80!black, line width=0.2pt, opacity=0.4] (6.560,0) to[out=58,in=122] (9.840,0);
  \draw[citewine, line width=0.2pt, opacity=0.4] (9.635,0) to[out=58,in=122] (10.045,0);
  \draw[projframe, line width=0.2pt, opacity=0.4] (8.610,0) to[out=58,in=122] (10.045,0);
  \draw[linknavy, line width=0.2pt, opacity=0.4] (6.970,0) to[out=58,in=122] (10.045,0);
  \draw[orange!80!black, line width=0.2pt, opacity=0.4] (6.765,0) to[out=58,in=122] (10.045,0);
  \draw[citewine, line width=0.2pt, opacity=0.4] (9.840,0) to[out=58,in=122] (10.250,0);
  \draw[projframe, line width=0.2pt, opacity=0.4] (8.815,0) to[out=58,in=122] (10.250,0);
  \draw[linknavy, line width=0.2pt, opacity=0.4] (7.175,0) to[out=58,in=122] (10.250,0);
  \draw[orange!80!black, line width=0.2pt, opacity=0.4] (6.970,0) to[out=58,in=122] (10.250,0);
  \draw[citewine, line width=0.2pt, opacity=0.4] (10.045,0) to[out=58,in=122] (10.455,0);
  \draw[projframe, line width=0.2pt, opacity=0.4] (9.020,0) to[out=58,in=122] (10.455,0);
  \draw[linknavy, line width=0.2pt, opacity=0.4] (7.380,0) to[out=58,in=122] (10.455,0);
  \draw[orange!80!black, line width=0.2pt, opacity=0.4] (7.175,0) to[out=58,in=122] (10.455,0);
  \draw[citewine, line width=0.2pt, opacity=0.4] (10.250,0) to[out=58,in=122] (10.660,0);
  \draw[projframe, line width=0.2pt, opacity=0.4] (9.225,0) to[out=58,in=122] (10.660,0);
  \draw[linknavy, line width=0.2pt, opacity=0.4] (7.585,0) to[out=58,in=122] (10.660,0);
  \draw[orange!80!black, line width=0.2pt, opacity=0.4] (7.380,0) to[out=58,in=122] (10.660,0);
  \draw[citewine, line width=0.2pt, opacity=0.4] (10.455,0) to[out=58,in=122] (10.865,0);
  \draw[projframe, line width=0.2pt, opacity=0.4] (9.430,0) to[out=58,in=122] (10.865,0);
  \draw[linknavy, line width=0.2pt, opacity=0.4] (7.790,0) to[out=58,in=122] (10.865,0);
  \draw[orange!80!black, line width=0.2pt, opacity=0.4] (7.585,0) to[out=58,in=122] (10.865,0);
  \draw[citewine, line width=0.2pt, opacity=0.4] (10.660,0) to[out=58,in=122] (11.070,0);
  \draw[projframe, line width=0.2pt, opacity=0.4] (9.635,0) to[out=58,in=122] (11.070,0);
  \draw[linknavy, line width=0.2pt, opacity=0.4] (7.995,0) to[out=58,in=122] (11.070,0);
  \draw[orange!80!black, line width=0.2pt, opacity=0.4] (7.790,0) to[out=58,in=122] (11.070,0);
  \draw[citewine, line width=0.2pt, opacity=0.4] (10.865,0) to[out=58,in=122] (11.275,0);
  \draw[projframe, line width=0.2pt, opacity=0.4] (9.840,0) to[out=58,in=122] (11.275,0);
  \draw[linknavy, line width=0.2pt, opacity=0.4] (8.200,0) to[out=58,in=122] (11.275,0);
  \draw[orange!80!black, line width=0.2pt, opacity=0.4] (7.995,0) to[out=58,in=122] (11.275,0);
  \draw[citewine, line width=0.2pt, opacity=0.4] (11.070,0) to[out=58,in=122] (11.480,0);
  \draw[projframe, line width=0.2pt, opacity=0.4] (10.045,0) to[out=58,in=122] (11.480,0);
  \draw[linknavy, line width=0.2pt, opacity=0.4] (8.405,0) to[out=58,in=122] (11.480,0);
  \draw[orange!80!black, line width=0.2pt, opacity=0.4] (8.200,0) to[out=58,in=122] (11.480,0);
  \draw[citewine, line width=0.2pt, opacity=0.4] (11.275,0) to[out=58,in=122] (11.685,0);
  \draw[projframe, line width=0.2pt, opacity=0.4] (10.250,0) to[out=58,in=122] (11.685,0);
  \draw[linknavy, line width=0.2pt, opacity=0.4] (8.610,0) to[out=58,in=122] (11.685,0);
  \draw[orange!80!black, line width=0.2pt, opacity=0.4] (8.405,0) to[out=58,in=122] (11.685,0);
  \draw[citewine, line width=0.2pt, opacity=0.4] (11.480,0) to[out=58,in=122] (11.890,0);
  \draw[projframe, line width=0.2pt, opacity=0.4] (10.455,0) to[out=58,in=122] (11.890,0);
  \draw[linknavy, line width=0.2pt, opacity=0.4] (8.815,0) to[out=58,in=122] (11.890,0);
  \draw[orange!80!black, line width=0.2pt, opacity=0.4] (8.610,0) to[out=58,in=122] (11.890,0);
  \draw[citewine, line width=0.2pt, opacity=0.4] (11.685,0) to[out=58,in=122] (12.095,0);
  \draw[projframe, line width=0.2pt, opacity=0.4] (10.660,0) to[out=58,in=122] (12.095,0);
  \draw[linknavy, line width=0.2pt, opacity=0.4] (9.020,0) to[out=58,in=122] (12.095,0);
  \draw[orange!80!black, line width=0.2pt, opacity=0.4] (8.815,0) to[out=58,in=122] (12.095,0);
  \draw[citewine, line width=0.2pt, opacity=0.4] (11.890,0) to[out=58,in=122] (12.300,0);
  \draw[projframe, line width=0.2pt, opacity=0.4] (10.865,0) to[out=58,in=122] (12.300,0);
  \draw[linknavy, line width=0.2pt, opacity=0.4] (9.225,0) to[out=58,in=122] (12.300,0);
  \draw[orange!80!black, line width=0.2pt, opacity=0.4] (9.020,0) to[out=58,in=122] (12.300,0);
  \draw[citewine, line width=0.2pt, opacity=0.4] (12.095,0) to[out=58,in=122] (12.505,0);
  \draw[projframe, line width=0.2pt, opacity=0.4] (11.070,0) to[out=58,in=122] (12.505,0);
  \draw[linknavy, line width=0.2pt, opacity=0.4] (9.430,0) to[out=58,in=122] (12.505,0);
  \draw[orange!80!black, line width=0.2pt, opacity=0.4] (9.225,0) to[out=58,in=122] (12.505,0);
  \draw[citewine, line width=0.2pt, opacity=0.4] (12.300,0) to[out=58,in=122] (12.710,0);
  \draw[projframe, line width=0.2pt, opacity=0.4] (11.275,0) to[out=58,in=122] (12.710,0);
  \draw[linknavy, line width=0.2pt, opacity=0.4] (9.635,0) to[out=58,in=122] (12.710,0);
  \draw[orange!80!black, line width=0.2pt, opacity=0.4] (9.430,0) to[out=58,in=122] (12.710,0);
  \draw[citewine, line width=0.2pt, opacity=0.4] (12.505,0) to[out=58,in=122] (12.915,0);
  \draw[projframe, line width=0.2pt, opacity=0.4] (11.480,0) to[out=58,in=122] (12.915,0);
  \draw[linknavy, line width=0.2pt, opacity=0.4] (9.840,0) to[out=58,in=122] (12.915,0);
  \draw[orange!80!black, line width=0.2pt, opacity=0.4] (9.635,0) to[out=58,in=122] (12.915,0);
  \fill[msgfill] (0.000,0) circle (1.5pt);
  \draw (0.000,0) circle (1.5pt);
  \fill[msgfill] (0.205,0) circle (1.5pt);
  \draw (0.205,0) circle (1.5pt);
  \fill[msgfill] (0.410,0) circle (1.5pt);
  \draw (0.410,0) circle (1.5pt);
  \fill[msgfill] (0.615,0) circle (1.5pt);
  \draw (0.615,0) circle (1.5pt);
  \fill[msgfill] (0.820,0) circle (1.5pt);
  \draw (0.820,0) circle (1.5pt);
  \fill[msgfill] (1.025,0) circle (1.5pt);
  \draw (1.025,0) circle (1.5pt);
  \fill[msgfill] (1.230,0) circle (1.5pt);
  \draw (1.230,0) circle (1.5pt);
  \fill[msgfill] (1.435,0) circle (1.5pt);
  \draw (1.435,0) circle (1.5pt);
  \fill[msgfill] (1.640,0) circle (1.5pt);
  \draw (1.640,0) circle (1.5pt);
  \fill[msgfill] (1.845,0) circle (1.5pt);
  \draw (1.845,0) circle (1.5pt);
  \fill[msgfill] (2.050,0) circle (1.5pt);
  \draw (2.050,0) circle (1.5pt);
  \fill[msgfill] (2.255,0) circle (1.5pt);
  \draw (2.255,0) circle (1.5pt);
  \fill[msgfill] (2.460,0) circle (1.5pt);
  \draw (2.460,0) circle (1.5pt);
  \fill[msgfill] (2.665,0) circle (1.5pt);
  \draw (2.665,0) circle (1.5pt);
  \fill[msgfill] (2.870,0) circle (1.5pt);
  \draw (2.870,0) circle (1.5pt);
  \fill[msgfill] (3.075,0) circle (1.5pt);
  \draw (3.075,0) circle (1.5pt);
  \fill[projframe!12] (3.280,0) circle (1.4pt);
  \fill[projframe!12] (3.485,0) circle (1.4pt);
  \fill[projframe!25] (3.690,0) circle (1.4pt);
  \fill[projframe!25] (3.895,0) circle (1.4pt);
  \fill[projframe!38] (4.100,0) circle (1.4pt);
  \fill[projframe!38] (4.305,0) circle (1.4pt);
  \fill[projframe!51] (4.510,0) circle (1.4pt);
  \fill[projframe!55] (4.715,0) circle (1.4pt);
  \fill[projframe!59] (4.920,0) circle (1.4pt);
  \fill[projframe!59] (5.125,0) circle (1.4pt);
  \fill[projframe!64] (5.330,0) circle (1.4pt);
  \fill[projframe!64] (5.535,0) circle (1.4pt);
  \fill[projframe!64] (5.740,0) circle (1.4pt);
  \fill[projframe!64] (5.945,0) circle (1.4pt);
  \fill[projframe!64] (6.150,0) circle (1.4pt);
  \fill[projframe!64] (6.355,0) circle (1.4pt);
  \fill[projframe!64] (6.560,0) circle (1.4pt);
  \fill[projframe!64] (6.765,0) circle (1.4pt);
  \fill[projframe!64] (6.970,0) circle (1.4pt);
  \fill[projframe!64] (7.175,0) circle (1.4pt);
  \fill[projframe!64] (7.380,0) circle (1.4pt);
  \fill[projframe!64] (7.585,0) circle (1.4pt);
  \fill[projframe!64] (7.790,0) circle (1.4pt);
  \fill[projframe!64] (7.995,0) circle (1.4pt);
  \fill[projframe!64] (8.200,0) circle (1.4pt);
  \fill[projframe!64] (8.405,0) circle (1.4pt);
  \fill[projframe!64] (8.610,0) circle (1.4pt);
  \fill[projframe!64] (8.815,0) circle (1.4pt);
  \fill[projframe!64] (9.020,0) circle (1.4pt);
  \fill[projframe!64] (9.225,0) circle (1.4pt);
  \fill[projframe!64] (9.430,0) circle (1.4pt);
  \fill[projframe!64] (9.635,0) circle (1.4pt);
  \fill[projframe!64] (9.840,0) circle (1.4pt);
  \fill[projframe!64] (10.045,0) circle (1.4pt);
  \fill[projframe!64] (10.250,0) circle (1.4pt);
  \fill[projframe!64] (10.455,0) circle (1.4pt);
  \fill[projframe!64] (10.660,0) circle (1.4pt);
  \fill[projframe!64] (10.865,0) circle (1.4pt);
  \fill[projframe!64] (11.070,0) circle (1.4pt);
  \fill[projframe!64] (11.275,0) circle (1.4pt);
  \fill[projframe!64] (11.480,0) circle (1.4pt);
  \fill[projframe!64] (11.685,0) circle (1.4pt);
  \fill[projframe!64] (11.890,0) circle (1.4pt);
  \fill[projframe!64] (12.095,0) circle (1.4pt);
  \fill[projframe!64] (12.300,0) circle (1.4pt);
  \fill[projframe!64] (12.505,0) circle (1.4pt);
  \fill[projframe!64] (12.710,0) circle (1.4pt);
  \fill[projframe!64] (12.915,0) circle (1.4pt);
  \draw[dashed, black!55] (5.330,-2.0) -- (5.330,-0.05);
  \node[black!70, anchor=west] at (5.410,-1.5) {\tiny $W_{26}$: all $16$ seeds reached};
  \draw[decorate, decoration={brace, mirror, amplitude=2.5pt}] (-0.03,-0.13) -- (3.075,-0.13) node[midway, below=1pt] {\tiny $16$ data seeds $W_0\dots W_{15}$};
  \draw[decorate, decoration={brace, mirror, amplitude=2.5pt}] (3.280,-0.13) -- (12.915,-0.13) node[midway, below=1pt] {\tiny $48$ derived $W_{16}\dots W_{63}$};
  \begin{scope}[yshift=-2.0cm]
    \draw[->, black!55] (-0.05,0) -- (13.120,0) node[right, black!70] {\tiny $t$};
    \draw[->, black!55] (-0.05,0) -- (-0.05,0.95) node[above, black!70, align=center] {\tiny seeds\\[-2pt]\tiny reached};
    \draw[black!30, dashed] (-0.05,0.8) -- (13.120,0.8) node[right, black!55] {\tiny $16$};
    \draw[projframe, thick] plot coordinates {(0.000,0.050) (0.205,0.050) (0.410,0.050) (0.615,0.050) (0.820,0.050) (1.025,0.050) (1.230,0.050) (1.435,0.050) (1.640,0.050) (1.845,0.050) (2.050,0.050) (2.255,0.050) (2.460,0.050) (2.665,0.050) (2.870,0.050) (3.075,0.050) (3.280,0.200) (3.485,0.200) (3.690,0.350) (3.895,0.350) (4.100,0.500) (4.305,0.500) (4.510,0.650) (4.715,0.700) (4.920,0.750) (5.125,0.750) (5.330,0.800) (5.535,0.800) (5.740,0.800) (5.945,0.800) (6.150,0.800) (6.355,0.800) (6.560,0.800) (6.765,0.800) (6.970,0.800) (7.175,0.800) (7.380,0.800) (7.585,0.800) (7.790,0.800) (7.995,0.800) (8.200,0.800) (8.405,0.800) (8.610,0.800) (8.815,0.800) (9.020,0.800) (9.225,0.800) (9.430,0.800) (9.635,0.800) (9.840,0.800) (10.045,0.800) (10.250,0.800) (10.455,0.800) (10.660,0.800) (10.865,0.800) (11.070,0.800) (11.275,0.800) (11.480,0.800) (11.685,0.800) (11.890,0.800) (12.095,0.800) (12.300,0.800) (12.505,0.800) (12.710,0.800) (12.915,0.800)};
  \end{scope}
\end{tikzpicture}
}
  \caption{The message schedule as a computed dependency graph
  (generated by \texttt{tools/wschedule\_graph.py}, which also emits a
  Graphviz \texttt{.dot} version). Nodes are the $64$ schedule words in
  index order; the sixteen data-carrying seeds $W_0,\dots,W_{15}$ are
  blue, and derived words darken with the number of seeds they
  transitively depend on. Each of the four arc colours is one term of
  $W_t=\sigma_1(W_{t-2})\shaplus W_{t-7}\shaplus\sigma_0(W_{t-15})\shaplus
  W_{t-16}$. The lower trace counts seeds reached versus $t$; it reaches
  all sixteen at $W_{26}$ (dashed line), after which the message is fully
  diffused through the schedule alone.}
  \label{fig:schedulegraph}
\end{figure}

One further structural reduction deserves emphasis, because it shrinks
the apparent dimensionality of the mystery. Only two of the eight
registers ever receive fresh values: each round computes one new word
entering the $a$-position and one entering the $e$-position, while the
other six merely shift along. The eight-word state is a \emph{delay
line}---a memoization of recent history, not eight independent degrees
of freedom. Writing $A_t$ and $E_t$ for the words produced at round $t$
(so that the state at that moment reads
$(A_t, A_{t-1}, A_{t-2}, A_{t-3},\, E_t, E_{t-1}, E_{t-2}, E_{t-3})$),
the sixty-four rounds are exactly the coupled two-track recurrence
\begin{align*}
  T_t &= E_{t-4} \shaplus \Sigma_1(E_{t-1})
        \shaplus \mathrm{Ch}(E_{t-1}, E_{t-2}, E_{t-3})
        \shaplus K_t \shaplus W_t,\\
  A_t &= T_t \shaplus \Sigma_0(A_{t-1})
        \shaplus \mathrm{Maj}(A_{t-1}, A_{t-2}, A_{t-3}),
  \qquad
  E_t \;=\; A_{t-4} \shaplus T_t,
\end{align*}
with the eight initial values $A_{-1}, \dots, A_{-4}, E_{-1}, \dots,
E_{-4}$ read off the incoming chaining state. So the ``$8$-dimensional''
iteration is really a pair of scalar recurrences over $\Ztwo$ with
lookback four---formally a nonlinear cousin of the lagged-Fibonacci and
linear-feedback recurrences whose theory is complete
(\cref{sec:prngs})---and the cryptanalytic literature does in fact work
in exactly this $A/E$ form. Nothing is lost in the reduction; the honest
shape of a round is two new words under an eight-word window of memory.

Notice what is \emph{absent}: no multiplications, no S-box tables, no
number-theoretic structure---only the three ARX primitives and two
quadratic Boolean functions, iterated $64$ times. The function is
degree-two nonlinearity compounded through $64$ rounds of geometric
whiplash between $\Ftwo^{256}$ and $(\Ztwo)^8$. That so spare a recipe
apparently produces a completely structureless map is the central mystery
this paper wants to sell.

\subsection{Two constructions that come with proofs}\label{sec:constructions}

It is worth pausing to locate the mystery precisely, because \SHA{} is
built in three nested layers and---remarkably---the two \emph{outer}
layers each carry a solid theorem. Only the innermost object is
theorem-free. Naming the layers from the outside in: the
\emph{mode} (Merkle--Damg{\aa}rd iteration, \cref{sec:merkledamgard})
wraps the \emph{compression function} $f$; the compression function is
itself built by the \emph{Davies--Meyer} construction around a block
cipher $E$; and $E$---the cipher SHACAL-2, which is just the round
machinery above run as a keyed permutation---is the finite map about
which nothing is proved. The provable security of \SHA{} is entirely a
statement about the two wrappers, conditional on the core behaving well.

\paragraph{The mode: Merkle--Damg{\aa}rd.} The outer layer's theorem is
already in hand: \cref{thm:md} says that a collision in the full hash
constructively yields a collision in $f$, so \emph{collision resistance
of $f$ is inherited by the whole function}---unconditionally, in the
standard model, needing only that the final block encode the message
length (the ``strengthening'' the proof consumes). This is as solid as
theorems in this subject get: no idealization, no unproven assumption, a
one-page reduction. What it does \emph{not} give is equally important and
sharply delimited---it transfers collision resistance and nothing else,
which is exactly why length extension and the multicollision family of
\cref{sec:mdstructure} slip through. The mode is a theorem about one
property, and an honest one.

\paragraph{The compression function: Davies--Meyer.} The middle layer is
the feedforward met in \cref{sec:roundshape}. Take a block cipher
$E_m(h)$---plaintext the chaining value $h\in\State$, key the message
block $m$---and set
\[
  f(h,m) \;=\; E_m(h) \,\shaplus\, h .
\]
The feedforward ``$\shaplus\,h$'' is the whole point. Without it,
$f = E_m$ would be a \emph{permutation} of $\State$ for each fixed $m$,
hence trivially invertible (decrypt with key $m$) and useless as a
one-way compression---the same observation that, read dynamically,
\cref{sec:unless} shows is what makes iterated \SHA{} lose entropy at
all. Adding the input back breaks invertibility, and does so in a way one
can prove is optimal, provided the cipher is idealized:

\begin{theorem}[Preneel--Govaerts--Vandewalle; Black--Rogaway--Shrimpton]
\label{thm:dm}
Model $E$ as an \emph{ideal cipher} (for each key an independent, uniformly
random permutation, accessible only as a black box in both directions).
Then the Davies--Meyer compression $f(h,m)=E_m(h)\shaplus h$ is optimally
collision- and preimage-resistant: any adversary asking $q$ cipher
queries finds a collision with probability at most $q(q+1)/2^{n}$, and a
preimage with probability of order $q/2^{n}$---matching the birthday
bound $2^{n/2}$ for collisions and the brute-force bound $2^{n}$ for
preimages, and both are tight.
\end{theorem}

\citet{PGV1993} isolated this by classifying all $64$ natural ways to
build a rate-one compression function from a block cipher and singling
out the $12$ sound ones (Davies--Meyer among them); \citet{BRS2002} then
supplied the rigorous black-box proofs and the exact bounds quoted above.
So the middle layer, too, has a theorem---but note the price on the
ticket. It holds in the \emph{ideal cipher model}, the block-cipher
cousin of the random-oracle model of \cref{sec:models}, and is therefore
a heuristic about \SHA{}, not a proof about it: SHACAL-2 is a fixed, public,
emphatically non-ideal cipher whose ``key schedule'' is literally the
message expansion of \cref{sec:roundshape}, so the assumption that
distinct message blocks act like independent random permutations is
exactly the leap of faith the model licenses and cannot justify. A
concrete reminder that the idealization is lossy: Davies--Meyer has
\emph{trivially computable fixed points}---solving $E_m(h)\shaplus h = h$
means $E_m(h)=0$, i.e.\ $h = E_m^{-1}(0)$, one short computation per
block---a harmless-looking feature that became a lever for the
long-message second-preimage attacks of \cref{sec:mdstructure}.

\paragraph{The moral, localized.} Two layers, two theorems: the mode
proved unconditionally, the compression proved in an idealized model, and
between them they reduce the security of the whole edifice to a single
demand---that the core map $f$ (equivalently the cipher $E$) have no
exploitable structure. That demand is \cref{sec:frameworks}'s gap, and
it is where every remaining page of this paper lives. The architecture is
solved mathematics; the mystery has been squeezed into one finite object,
and the constructions of this section are precisely the machinery that
did the squeezing.

\subsection{One honest theorem: invertible diffusion}\label{sec:rivest}

Amid so much that is unproved, it is worth exhibiting in full one of the
few genuine theorems about \SHA{}'s internals---partly because it is
lovely, partly because its proof lives in exactly the finite geometry of
\cref{sec:statespace}, and partly because it shows the ``middle
territory'' of \cref{sec:gap} is not empty by necessity.

Ask a basic anatomical question about the stirring maps of
\cref{sec:roundshape}: are $\Sigma_0, \Sigma_1, \sigma_0, \sigma_1$
\emph{bijections} of the word space $\Ftwo^{32}$? (They are
$\Ftwo$-linear, hence ``nonlinearities'' only through the other lens of
\cref{sec:statespace}---as maps of $\Ztwo$ they are wildly nonlinear.
The question is whether they discard information.) Rivest answered this
for the pure-rotation case in a four-page note~\citep{Rivest2011}:

\begin{theorem}[Rivest]\label{thm:rivest}
Let $w$ be a power of two. The exclusive-or of $t$ distinct rotations of
a $w$-bit word,
$x \mapsto \mathrm{ROTR}^{r_1}(x) \shaxor \cdots \shaxor \mathrm{ROTR}^{r_t}(x)$,
is invertible if and only if $t$ is odd.
\end{theorem}

\begin{proof}[Proof idea]
Identify $\Ftwo^{32}$ with the quotient ring
$R = \Ftwo[x]/(x^{32}-1)$, a bit vector becoming a polynomial and
rotation becoming multiplication by $x^{r}$. The map in question is
multiplication by $p(x) = x^{r_1} + \cdots + x^{r_t}$, invertible in $R$
exactly when $\gcd\bigl(p(x),\, x^{32}-1\bigr) = 1$. Over $\Ftwo$ the
modulus collapses---$x^{32} - 1 = (x-1)^{32}$, because $32$ is a power of
two and squaring is a ring homomorphism in characteristic $2$---so the
gcd is $1$ precisely when $(x-1) \nmid p$, i.e.\ when $p(1) = t \bmod 2$
is $1$.
\end{proof}

So $\Sigma_0 = \mathrm{ROTR}^{2} \shaxor \mathrm{ROTR}^{13} \shaxor
\mathrm{ROTR}^{22}$ and $\Sigma_1$, with their three rotations each, are
bijections: they stir without spilling. The theorem repays contemplation.
It is two lines of algebra, yet it turns on the arithmetic accident that
the word length is a power of two, which makes the cyclic algebra $R$
\emph{local} (one maximal ideal, generated by $x-1$): the only way an
xor-of-rotations can fail to be invertible is by having an even number of
terms. Had the designers used $33$-bit words, $x^{33}-1$ would factor
richly and the parity rule would be false. Small finite geometries have
loud opinions.

Because the proof is constructive, invertibility need not be left as an
existence statement: the inverses can be exhibited outright. Carrying out
the division in $R$---equivalently, inverting the corresponding
$32\times32$ circulant matrices over $\Ftwo$; \cref{app:sigmacode} gives
a complete reference implementation, which also confirms the results on
thousands of random words---produces closed forms that are again xors of
rotations:
\begin{equation}\label{eq:sigmainv}
\begin{aligned}
  \Sigma_0^{-1} \;&=\; \bigoplus_{r \in S_0} \mathrm{ROTR}^{r},
  & S_0 &= \{0,1,2,5,7,8,9,10,12,16,17,19,22,23,25,29,30\},\\
  \Sigma_1^{-1} \;&=\; \bigoplus_{r \in S_1} \mathrm{ROTR}^{r},
  & S_1 &= \{2,3,5,7,9,11,12,13,18,21,22,23,24,26,27,29,30\},
\end{aligned}
\end{equation}
each an xor of \emph{seventeen} rotations ($S_0$ includes $r = 0$, the
identity rotation). Seventeen is odd, as \cref{thm:rivest} demands of any
invertible xor-of-rotations---and as it must be on general grounds:
evaluating $p\,q \equiv 1$ at $x = 1$ forces
$|p| \cdot |q| \equiv 1 \pmod 2$, so odd-weight elements have odd-weight
inverses. The asymmetry of sparseness is the striking part: three
rotations to stir, seventeen to unstir---more than half of the
thirty-two possible offsets. Invertible is not the same as symmetrically
cheap, even inside the linear layer.

The ring picture yields more than the inverses; it computes the whole
multiplicative life of these maps, essentially by hand. Composition of
rotations adds offsets modulo $32$, so write an xor-of-rotations as a
polynomial with the \emph{offsets as exponents}: $\Sigma_0$ becomes
$p_0 = x^{2} + x^{13} + x^{22}$. Squaring is the Frobenius map of
characteristic $2$---it simply doubles every offset---so squaring $p_0$
three times gives, after reducing exponents modulo $32$,
$\Sigma_0^{8} = \mathrm{ROTR}^{8}$: \emph{eight applications of the
three-rotation stir collapse to one bare rotation}. Hence
$\Sigma_0^{32} = \mathrm{id}$, and the order of $\Sigma_0$ in
$\mathrm{GL}_{32}(\Ftwo)$ is exactly $32$; the same computation gives
$\Sigma_1^{16} = \mathrm{id}$ with order exactly $16$. In particular
$\Sigma_0^{-1} = \Sigma_0^{31}$ and $\Sigma_1^{-1} = \Sigma_1^{15}$, so
\cref{eq:sigmainv} could have been generated, offsets and all, by
repeated squaring. One more turn of the crank and the observation
becomes a theorem: for \emph{any} odd xor-of-rotations $u$ of $32$-bit
words, $u^{32} = u(x)^{32} = u(x^{32}) = u(1) = 1$, so every invertible
xor-of-rotations has order dividing $32$. The stirring maps of \SHA{},
for all their visual violence, are elements of small $2$-power order in
a $2$-group---one more exact structural fact living quietly next door to
the conjectured structurelessness (\cref{sec:hope}).

Note what the theorem does \emph{not} cover: $\sigma_0 =
\mathrm{ROTR}^{7} \shaxor \mathrm{ROTR}^{18} \shaxor \mathrm{SHR}^{3}$
and $\sigma_1$ each replace one rotation by a bit-discarding shift,
leaving the polynomial ring for plain linear algebra; and the genuinely
non-invertible ingredients of \SHA{} are elsewhere entirely---the
bitwise-quadratic $\mathrm{Ch}$ and $\mathrm{Maj}$ and the Davies--Meyer
feedforward, where all the designed information loss is concentrated. The
picture that emerges is architectural: \emph{lossless} stirring layers
(provably so, by \cref{thm:rivest}), \emph{lossy} mixing concentrated in
three auditable places. That so clean a structural statement was worth a
research note in 2011---twenty years into the design lineage---says
something about how sparsely theorized this territory is.

\begin{project}{The linear algebra of diffusion}{linear algebra;
  programming}{3--4 weeks}
  Write each of $\Sigma_0, \Sigma_1, \sigma_0, \sigma_1$ as an explicit
  $32 \times 32$ matrix over $\Ftwo$ and decide invertibility of the
  $\sigma$'s (settling what \cref{thm:rivest} leaves open). Compute the
  order of each invertible map in $\mathrm{GL}_{32}(\Ftwo)$, the cycle
  structure of its action on $\Ftwo^{32}$, and the subgroup the four maps
  generate. Then reprove \cref{thm:rivest} for general word length $w$:
  the answer is a pretty gcd condition (worked out in Rivest's note), and
  comparing $w = 32$ with $w = 33$ makes the power-of-two accident vivid.
\end{project}

\begin{project}{Deriving the exact inverses of $\Sigma_0$ and $\Sigma_1$}
  {polynomial arithmetic over $\Ftwo$; programming}{2--3 weeks}
  \Cref{thm:rivest} guarantees that $\Sigma_0$ and $\Sigma_1$ are
  bijections, and \cref{eq:sigmainv} displays their inverses; this
  project derives that display from scratch, by two independent methods,
  and decides what \cref{eq:sigmainv} deliberately leaves out. The steps
  are fully determined:
  \begin{enumerate}
    \item \textbf{Move to the ring.} Identify a word with a polynomial in
      $R = \Ftwo[x]/(x^{32}-1)$, bit $i$ becoming the coefficient of
      $x^i$, so that a rotation is multiplication by a power of $x$.
      Write $\Sigma_0$ as multiplication by
      $p_0(x) = x^{2} + x^{13} + x^{22}$ and $\Sigma_1$ by
      $p_1(x) = x^{6} + x^{11} + x^{25}$ (the \SHA{} rotation offsets).
    \item \textbf{Invert by linear algebra.} Multiplication by $p$ is an
      $\Ftwo$-linear circulant map; build its $32\times32$ matrix (columns
      are $p, xp, x^2p, \dots$) and invert it over $\Ftwo$ by Gaussian
      elimination. The last column of the inverse, read as a polynomial
      $q$, satisfies $p\,q \equiv 1 \pmod{x^{32}-1}$.
    \item \textbf{Read off the answer.} You should recover exactly the
      seventeen-term rotation sets of \cref{eq:sigmainv}. Verify
      $p\,q = 1$ in $R$ directly, and check your derivation against the
      reference implementation of \cref{app:sigmacode}. Then reproduce
      \cref{eq:sigmainv} a third way, with no linear algebra at all: by
      repeated squaring, using $\Sigma_0^{-1} = \Sigma_0^{31}$ and
      $\Sigma_1^{-1} = \Sigma_1^{15}$.
    \item \textbf{Cross-check by Euclid.} Recompute $q$ by the extended
      Euclidean algorithm on $p(x)$ and $x^{32}-1 = (x-1)^{32}$ in
      $\Ftwo[x]$; the two methods must agree.
    \item \textbf{Contrast with the $\sigma$'s.} Repeat for
      $\sigma_0 = \mathrm{ROTR}^{7}\shaxor\mathrm{ROTR}^{18}\shaxor
      \mathrm{SHR}^{3}$: the shift $\mathrm{SHR}^{3}$ leaves the cyclic
      ring (it is not a unit multiple of $x$), so the polynomial trick
      fails and only the matrix method applies. Decide invertibility of
      $\sigma_0, \sigma_1$ this way---settling what \cref{thm:rivest}
      does not cover.
  \end{enumerate}
  Deliverable: the two explicit inverse maps, a proof they are correct,
  and a one-paragraph account of why the $\sigma$'s need the heavier tool.
\end{project}

\subsection{Nothing up my sleeve}\label{sec:nums}

One question a mathematician should ask of \cref{sec:roundshape}: where do
the constants come from? The initial state $\mathrm{IV}$ and the sixty-four
round constants $K_t$ could, in principle, hide a trapdoor---constants
engineered so their designer knows structure nobody else can see. The
history is not hypothetical: the unexplained S-box constants of the 1970s
DES cipher fed decades of suspicion (later resolved in the designers'
favor), and post-Snowden, an actual back-doored NIST standard (the
Dual\_EC random generator) was withdrawn. The SHA family itself supplies
the sharpest demonstration that its designer's choices are informed
rather than arbitrary: the NSA's unexplained one-bit revision of SHA-0
into SHA-1 (see the footnote in \cref{sec:history}) anticipated by four
years the differential attack later found in the open
literature~\citep{ChabaudJoux1998}. The SHA-2 designers therefore
committed to constants they visibly could not have chosen:
\[
  \mathrm{IV}_j = \text{first 32 fractional bits of } \sqrt{p_j},
  \qquad
  K_t = \text{first 32 fractional bits of } \sqrt[3]{p_t},
\]
where $p_1 < p_2 < \cdots$ are the primes $2, 3, 5, \dots$~\citep{FIPS180-4}.
(So $\mathrm{IV}_1 = \mathtt{6a09e667}$ encodes $\sqrt{2}-1$ in binary.)
Such choices are called \emph{nothing-up-my-sleeve} constants: points of
the state space selected by a rule so canonical that no freedom remains in
which to hide intent.

There is a pleasing irony here that deserves a mathematician's smile: the
constants are trusted \emph{because} they are digits of irrational
numbers, yet essentially nothing is proved about those digits either---the
normality of $\sqrt{2}$ is an open problem. The construction does not
replace an unproved assumption with a proved one; it replaces an
\emph{unauditable} choice with an \emph{auditable} one. Epistemology, not
theorems, is doing the work---a theme by now familiar
(\cref{sec:epistemology}).

\subsection{Reasons for hope: exact solutions next door}\label{sec:hope}

Almost everything after this section will be statistical:
random-mapping silhouettes, test batteries, mixing arguments---the
approximate machinery one reaches for when exact structure is
unavailable. Before conceding that this is the only possible register,
it is worth assembling the evidence that exact structure keeps turning
up \emph{next door}, because the history of discrete iteration is
littered with problems that looked hopelessly procedural until someone
found the right coordinates.

The Josephus process (\cref{sec:josephus}) is this paper's opening
specimen: an $(n-1)$-step elimination dynamics collapsed to a single
bit rotation. \Cref{thm:rivest} is a second, living inside \SHA{}
itself: a question about its diffusion maps settled by two lines of
algebra in a local ring. The genre is much larger. Additive cellular
automata---iterative, visually turbulent, every bit as ``chaotic'' in
appearance as a hash round---were solved outright by
\citet{MartinOdlyzkoWolfram1984}: represent configurations as
polynomials over $\Ftwo$ and time evolution becomes multiplication,
with global transient and cycle structure falling out as exact
divisibility statements. \citet{Sutner1989} dispatched the Lights-Out
puzzle the same way (the $\sigma$-game as $\Ftwo$-linear algebra). The
pattern is older than computing: the logistic map at parameter $4$, the
textbook ``chaotic'' system, has the closed form
$x_t = \sin^2\!\bigl(2^t \arcsin\sqrt{x_0}\,\bigr)$---an observation
going back to Ulam and von Neumann, and the simplest instance of
Schr{\"o}der's program of solving iteration by
conjugation~\citep{Schroeder1871}: find $\psi$ making
$\psi \circ f \circ \psi^{-1}$ linear, and the $t$-th iterate is
computed in one stroke. That program grew into modern complex
dynamics~\citep{Milnor2006}. ``Chaotic'' and ``exactly solvable'' are
not antonyms; they are frequently the same object described before and
after the right change of variables.

Within the hash family the lesson recurs. SHA-1's message schedule is
$\Ftwo$-linear, hence \emph{exactly} analyzable---its difference
propagation is linear algebra---and that handle is precisely what the
collision attacks of \cref{sec:falling} gripped. The professional
response to such objects has been the cryptanalyst's panoply,
differential attacks in the tradition of \citet{BihamShamir1991} and
their descendants: a \emph{statistical} theory of near-structure,
probability-weighted trails threaded through the rounds
(\cref{sec:differential} is devoted to it). The examples
above argue for keeping a different question alive: whether some change
of coordinates---a direct finite-geometry path, in the spirit of the
polynomial representations that solved cellular automata, with no
mixing theory, no ergodicity, and nothing approximate---renders some
nontrivial slice of the iteration exactly computable. The designers'
explicit goal was to foreclose this (\cref{sec:statespace}), and the
smart money says they succeeded. But the foreclosure is a conjecture,
not a theorem; the components keep falling one at a time to exact
methods (\cref{thm:rivest}, the T-function
theory~\citep{KlimovShamir2003, Anashin2006}, the two-track reduction
of \cref{sec:roundshape}); and the hunting grounds are tiny finite
objects, where undergraduate computation is a genuine research
instrument (\cref{sec:projects}).
